\documentclass[12pt,a4paper]{article}
\usepackage[utf8]{vietnam}
\usepackage{amsmath,amssymb}
\usepackage{geometry}
\geometry{top=1.5cm,bottom=1.2cm,left=1.4cm,right=1.4cm,headheight=28pt,headsep=12pt,footskip=18pt}
\usepackage{tikz}
\usepackage{xcolor}
\usepackage{enumerate}
\usepackage{graphicx}
\usepackage[version=4]{mhchem}
\usepackage{fancyhdr}
\usepackage{ex_test}

\renewcommand{\baselinestretch}{1.0}
\setlength{\parskip}{1pt}

\definecolor{hfRed}{RGB}{211,47,47}
\definecolor{hfGreen}{RGB}{139,195,144}

\pagestyle{fancy}
\fancyhf{}
\fancyhead[L]{\small\itshape\color{hfRed} Tài liệu lưu hành nội bộ}
\fancyhead[R]{\small\itshape\color{hfRed} Thầy \textbf{TRẦN VĂN THẠNH} - 0777.470.803}
\renewcommand{\headrulewidth}{0.8pt}
\renewcommand{\headrule}{\hbox to\headwidth{\color{hfGreen}\leaders\hrule height \headrulewidth\hfill}}

\fancyfoot[L]{\small\itshape\color{hfRed} CS1: 6/15 Nguyễn Hoàng, P. Kim Long, TP Huế \quad|\quad CS2: 24 Đặng Thái Thân, TP Huế}
\fancyfoot[R]{\small\bfseries\color{hfRed} Trang \thepage}
\renewcommand{\footrulewidth}{0.8pt}
\renewcommand{\footrule}{\hbox to\headwidth{\color{hfGreen}\leaders\hrule height \footrulewidth\hfill}}

\fancypagestyle{firstpage}{
  \fancyhf{}
  \renewcommand{\headrulewidth}{0pt}
  \fancyfoot[L]{\small\itshape\color{hfRed} CS1: 6/15 Nguyễn Hoàng, P. Kim Long, TP Huế \quad|\quad CS2: 24 Đặng Thái Thân, TP Huế}
  \fancyfoot[R]{\small\bfseries\color{hfRed} Trang \thepage}
  \renewcommand{\footrulewidth}{0.8pt}
  \renewcommand{\footrule}{\hbox to\headwidth{\color{hfGreen}\leaders\hrule height \footrulewidth\hfill}}
}

\begin{document}
\thispagestyle{firstpage}

\noindent
\begin{minipage}[t]{0.45\textwidth}
	\centering\fontsize{10pt}{12pt}\selectfont
	\textbf{SỞ GIÁO DỤC VÀ ĐÀO TẠO THÀNH PHỐ HUẾ}\\
	\textbf{TRƯỜNG THPT HOÁ CHÂU}\\
	\textbf{ĐỀ CHÍNH THỨC}\\
	\textit{(Đề thi có 03 trang)}
\end{minipage}\hfill
\begin{minipage}[t]{0.53\textwidth}
	\centering\small
	\textbf{KIỂM TRA GIỮA HỌC KỲ I -- NĂM HỌC 2025 - 2026}\\
	\textbf{MÔN: HOÁ HỌC -- LỚP 11}\\
	\textit{Thời gian làm bài: 45 phút (không kể thời gian phát đề)}\\
	\fbox{\textbf{MÃ ĐỀ: 102}}
\end{minipage}

\smallskip
\begin{center}
	{\color{hfRed}\textbf{\large (BẢN GIÁO VIÊN -- CÓ LỜI GIẢI CHI TIẾT)}}
\end{center}

\smallskip
\noindent\hrulefill
\smallskip

\noindent\textbf{Họ, tên thí sinh:} \dotfill\ \textbf{Lớp:} \dotfill\ \textbf{Số báo danh:} \dotfill

\smallskip
\noindent\textit{(Cho biết nguyên tử khối: H = 1; C = 12; N = 14; O = 16; Na = 23; Al = 27; S = 32; Cl = 35,5; K = 39; Ca = 40; Fe = 56; Cu = 64; Ba = 137).}
\smallskip

\noindent
\textbf{\large A. PHẦN TRẮC NGHIỆM (7,0 điểm)}

\smallskip
\noindent
\textbf{PHẦN I (3,0 điểm): Câu trắc nghiệm nhiều phương án lựa chọn.}\\
\textit{Thí sinh trả lời từ câu 1 đến câu 12. Mỗi câu hỏi thí sinh chỉ chọn một phương án.}

\Opensolutionfile{ans}[ans-hoa11-ma102gv]

\begin{ex}
Cho phương trình: $\ce{CH3COOH + H2O <=> CH3COO- + H3O+}$. Trong phản ứng thuận, theo thuyết Brønsted -- Lowry chất nào là acid?
\choice
{\True $\ce{CH3COOH}$.}
{$\ce{CH3COO-}$.}
{$\ce{H3O+}$.}
{$\ce{H2O}$.}
\loigiai{
Theo thuyết Brønsted -- Lowry, acid là chất cho proton ($\ce{H+}$). Trong phản ứng thuận, $\ce{CH3COOH}$ nhường proton $\ce{H+}$ cho $\ce{H2O}$ để tạo thành $\ce{CH3COO-}$, do đó $\ce{CH3COOH}$ là acid.\\
\textbf{Chọn A.}}
\end{ex}

\begin{ex}
Vai trò của $\ce{NH3}$ trong phản ứng: $\ce{4NH3 + 5O2 ->[t^\circ, \text{Pt}] 4NO + 6H2O}$ là
\choice
{\True chất khử.}
{chất oxi hóa.}
{acid.}
{base.}
\loigiai{
Số oxi hóa của N trong $\ce{NH3}$ là $-3$, sau phản ứng trong $\ce{NO}$ là $+2$ (tăng từ $-3$ lên $+2$). Vì số oxi hóa tăng nên $\ce{NH3}$ nhường electron, đóng vai trò là chất khử.\\
\textbf{Chọn A.}}
\end{ex}

\begin{ex}
Để xác định nồng độ của một dung dịch $\ce{NaOH}$, người ta đã tiến hành chuẩn độ bằng dung dịch $\ce{HCl}\ 0{,}1\text{ M}$. Để chuẩn độ $10\text{ mL}$ dung dịch $\ce{NaOH}$ này cần $12\text{ mL}$ dung dịch $\ce{HCl}$. Nồng độ của dung dịch $\ce{NaOH}$ trên là
\choice
{$1{,}2\text{ M}$.}
{$0{,}012\text{ M}$.}
{$0{,}1\text{ M}$.}
{\True $0{,}12\text{ M}$.}
\loigiai{
Phương trình phản ứng chuẩn độ: $\ce{NaOH + HCl -> NaCl + H2O}$.\\
Tại điểm tương đương:
\[ n_{\ce{NaOH}} = n_{\ce{HCl}} \iff C_{\ce{NaOH}} \times V_{\ce{NaOH}} = C_{\ce{HCl}} \times V_{\ce{HCl}} \]
\[ \implies C_{\ce{NaOH}} = \frac{0{,}1 \times 12}{10} = 0{,}12\text{ M}. \]
\textbf{Chọn D.}}
\end{ex}

\begin{ex}
Phản ứng thuận nghịch là phản ứng
\choice
{xảy ra giữa hai chất khí.}
{\True xảy ra theo hai chiều ngược nhau trong cùng điều kiện.}
{có phương trình hoá học được biểu diễn bằng mũi tên một chiều.}
{chỉ xảy ra theo một chiều nhất định.}
\loigiai{
Phản ứng thuận nghịch là phản ứng xảy ra theo hai chiều ngược nhau trong cùng điều kiện xác định.\\
\textbf{Chọn B.}}
\end{ex}

\begin{ex}
$\ce{HNO3}$ tinh khiết là chất lỏng không màu, nhưng dung dịch $\ce{HNO3}$ để lâu thường chuyển sang màu vàng là do
\choice
{khi để lâu $\ce{HNO3}$ bị khử bởi các chất của môi trường.}
{$\ce{HNO3}$ tan nhiều trong nước.}
{dung dịch $\ce{HNO3}$ có tính oxi hóa mạnh.}
{\True dung dịch $\ce{HNO3}$ có hoà tan một lượng nhỏ $\ce{NO2}$.}
\loigiai{
Dưới tác dụng của ánh sáng hoặc nhiệt độ, $\ce{HNO3}$ bị phân hủy một phần:\\
$\ce{4HNO3 -> 4NO2 + O2 + 2H2O}$.\\
Khí $\ce{NO2}$ sinh ra có màu nâu đỏ tan lại trong dung dịch acid làm cho dung dịch chuyển sang màu vàng.\\
\textbf{Chọn D.}}
\end{ex}

\begin{ex}
Cho cân bằng (trong bình kín) sau:
\[ \ce{CO(g) + H2O(g) <=> CO2(g) + H2(g)} \quad \Delta_r H_{298}^\circ < 0 \]
Trong các yếu tố: (1) tăng nhiệt độ; (2) thêm một lượng hơi nước; (3) thêm một lượng $\ce{H2}$; (4) tăng áp suất chung của hệ; (5) dùng chất xúc tác. Dãy gồm các yếu tố đều làm thay đổi cân bằng của hệ là
\choice
{(2), (3), (4).}
{\True (1), (2), (3).}
{(1), (4), (5).}
{(1), (2), (4).}
\loigiai{
- (1) Tăng nhiệt độ: do phản ứng tỏa nhiệt ($\Delta_r H_{298}^\circ < 0$), khi tăng nhiệt độ cân bằng chuyển dịch theo chiều thu nhiệt (chiều nghịch).\\
- (2) Thêm hơi nước: tăng nồng độ chất phản ứng, cân bằng chuyển dịch theo chiều thuận.\\
- (3) Thêm $\ce{H2}$: tăng nồng độ sản phẩm, cân bằng chuyển dịch theo chiều nghịch.\\
- (4) Áp suất: do tổng số mol khí hai vế bằng nhau ($1 + 1 = 1 + 1 = 2$) nên áp suất không làm chuyển dịch cân bằng.\\
- (5) Chất xúc tác: chỉ làm tăng tốc độ phản ứng chứ không làm chuyển dịch cân bằng.\\
Vậy dãy gồm các yếu tố đều làm thay đổi cân bằng là (1), (2), (3).\\
\textbf{Chọn B.}}
\end{ex}

\begin{ex}
Phát biểu nào dưới đây \textbf{không đúng}?
\choice
{$\ce{H2SO4}$ đặc là chất hút nước mạnh.}
{\True Khi pha loãng sulfuric acid chỉ được cho từ từ nước vào acid.}
{$\ce{H2SO4}$ loãng có đầy đủ tính chất chung của acid.}
{Khi tiếp xúc với $\ce{H2SO4}$ đặc dễ gây bỏng nặng.}
\loigiai{
Quy tắc pha loãng sulfuric acid an toàn: rót từ từ acid vào nước và khuấy nhẹ bằng đũa thủy tinh, tuyệt đối không làm ngược lại (không rót nước vào acid) vì acid sulfuric đặc tỏa nhiệt rất lớn làm nước sôi đột ngột và bắn các giọt acid nóng ra ngoài gây bỏng nguy hiểm. Do đó phát biểu B là không đúng.\\
\textbf{Chọn B.}}
\end{ex}

\begin{ex}
Cho các phát biểu sau:\\
(a) Sulfur là chất rắn, màu trắng, không tan trong nước.\\
(b) Sulfur và sulfur dioxide vừa có tính oxi hóa, vừa có tính khử.\\
(c) Khi phản ứng với hydrogen, sulfur thể hiện tính oxi hóa.\\
(d) Nước thải sinh hoạt là một trong các nguồn phát thải khí $\ce{SO2}$.\\
(e) Sulfur dioxide được sử dụng để tẩy trắng vải sợi, bột giấy, sản xuất sulfuric acid và diệt nấm mốc.\\
Số phát biểu \textbf{đúng} là
\choice
{4.}
{1.}
{2.}
{\True 3.}
\loigiai{
- (a) Sai vì sulfur là chất rắn màu vàng (không phải màu trắng).\\
- (b) Đúng vì S(0) và S(+4) trong $\ce{SO2}$ có số oxi hóa trung gian nên vừa có tính oxi hóa vừa có tính khử.\\
- (c) Đúng vì trong phản ứng $\ce{H2 + S -> H2S}$, số oxi hóa của S giảm từ 0 xuống $-2$ nên S là chất oxi hóa.\\
- (d) Sai vì nguồn phát thải $\ce{SO2}$ chủ yếu từ đốt than đá, dầu mỏ, luyện kim màu và khí núi lửa; nước thải sinh hoạt chủ yếu sinh ra $\ce{H2S}$, $\ce{NH3}$, $\ce{CH4}$.\\
- (e) Đúng vì $\ce{SO2}$ được dùng để tẩy trắng vải sợi, bột giấy, sản xuất sulfuric acid và diệt nấm mốc.\\
Vậy có 3 phát biểu đúng: (b), (c), (e).\\
\textbf{Chọn D.}}
\end{ex}

\begin{ex}
Để tạo độ xốp cho một số loại bánh, có thể dùng chất nào sau đây?
\choice
{\True $\ce{NH4HCO3}$.}
{$\ce{CaCO3}$.}
{$\ce{NaCl}$.}
{$\ce{(NH4)3PO4}$.}
\loigiai{
$\ce{NH4HCO3}$ (bột nở) khi nướng bánh bị nhiệt phân tạo ra các chất khí $\ce{NH3}$, $\ce{CO2}$ và hơi nước thoát ra tạo các lỗ xốp li ti trong bánh: $\ce{NH4HCO3 ->[t^\circ] NH3 ^ + CO2 ^ + H2O ^}$.\\
\textbf{Chọn A.}}
\end{ex}

\begin{ex}
Ở nhiệt độ thường, nitrogen khá trơ về mặt hoạt động là do
\choice
{\True Phân tử nitrogen có liên kết ba khá bền.}
{Nitrogen có độ âm điện lớn nhất trong nhóm.}
{Phân tử nitrogen không phân cực.}
{Nitrogen có bán kính nguyên tử nhỏ.}
\loigiai{
Phân tử nitrogen gồm 2 nguyên tử liên kết với nhau bằng một liên kết ba ($\ce{N#N}$) rất bền với năng lượng liên kết rất lớn ($945\text{ kJ/mol}$), do đó ở điều kiện thường nitrogen trơ về mặt hóa học.\\
\textbf{Chọn A.}}
\end{ex}

\begin{ex}
Tính chất nào sau đây \textbf{không phải} tính chất vật lí của sulfur?
\choice
{Thể rắn ở điều kiện thường.}
{Không tan trong nước.}
{Màu vàng ở điều kiện thường.}
{\True Không tan trong benzene.}
\loigiai{
Sulfur không tan trong nước nhưng tan nhiều trong các dung môi hữu cơ không phân cực như benzene ($\ce{C6H6}$), toluene, carbon disulfide ($\ce{CS2}$)... Do đó "Không tan trong benzene" là phát biểu không đúng.\\
\textbf{Chọn D.}}
\end{ex}

\begin{ex}
Nitric acid thường được sử dụng để tẩy rửa ống, bề mặt kim loại do có
\choice
{tính khử.}
{\True tính oxi hoá mạnh.}
{tính acid yếu.}
{tính khử và tính acid mạnh.}
\loigiai{
Nitric acid có tính oxi hóa rất mạnh và tính acid mạnh, có khả năng hòa tan các oxit kim loại cũng như hầu hết các kim loại bám bẩn trên bề mặt đường ống để làm sạch kim loại.\\
\textbf{Chọn B.}}
\end{ex}

\Closesolutionfile{ans}

\bigskip
\noindent
\textbf{PHẦN II (2,0 điểm): Câu trắc nghiệm đúng sai.}\\
\textit{Thí sinh trả lời từ câu 1 đến câu 2. Trong mỗi ý a), b), c), d) ở mỗi câu, thí sinh chọn đúng hoặc sai.}

\begin{ex}
Cho phản ứng tổng hợp ammonia:
\[ \ce{N2(g) + 3H2(g) <=>[xt, t^\circ, p] 2NH3(g)} \quad \Delta_r H_{298}^\circ = -92\text{ kJ} \]
\choiceTF
{\True a) Khi tăng nhiệt độ của hệ thì cân bằng chuyển dịch theo chiều nghịch.}
{b) Phản ứng trên là phản ứng thu nhiệt theo chiều thuận.}
{c) Để phản ứng tổng hợp $\ce{NH3}$ đạt hiệu suất cao thì người ta thực hiện ở nhiệt độ rất thấp và áp suất cao.}
{d) Trong phản ứng trên $\ce{N2}$ đóng vai trò là chất khử.}
\loigiai{
a) Đúng vì phản ứng thuận có $\Delta_r H_{298}^\circ = -92\text{ kJ} < 0$ (phản ứng tỏa nhiệt). Khi tăng nhiệt độ, cân bằng chuyển dịch theo chiều thu nhiệt (chiều nghịch).\\
b) Sai vì $\Delta_r H_{298}^\circ < 0$ nên phản ứng thuận là phản ứng tỏa nhiệt.\\
c) Sai vì ở nhiệt độ rất thấp tốc độ phản ứng xảy ra vô cùng chậm và chất xúc tác không hoạt động. Trong công nghiệp, phản ứng được thực hiện ở nhiệt độ thích hợp (khoảng $400 - 450^\circ\text{C}$) và áp suất cao ($150 - 200\text{ bar}$) có mặt xúc tác sắt.\\
d) Sai vì số oxi hóa của N giảm từ 0 ($\ce{N2}$) xuống $-3$ ($\ce{NH3}$) nên $\ce{N2}$ đóng vai trò là chất oxi hóa.}
\end{ex}

\begin{ex}
Chọn phát biểu đúng -- sai:
\choiceTF
{a) Dung dịch sulfuric acid đặc hòa tan được tất cả các kim loại.}
{\True b) Hòa tan $18{,}4\text{ gam}$ hỗn hợp X gồm Fe và Cu trong dung dịch $\ce{H2SO4}$ (đặc, nóng, dư), thu được $8{,}6765\text{ Lít}$ $\ce{SO2}$ (đkc) và dung dịch Y. Phần trăm khối lượng của Fe trong hỗn hợp X là $30{,}43\%$.}
{c) Dung dịch $\ce{HNO3}$ đặc, nguội phản ứng được với Al và Fe.}
{\True d) Sulfuric acid tan tốt trong nước, quá trình hòa tan tỏa nhiệt mạnh.}
\loigiai{
a) Sai vì $\ce{H2SO4}$ đặc không hòa tan được vàng (Au) và platinum (Pt).\\
b) Đúng. Số mol $\ce{SO2} = \frac{8{,}6765}{24{,}79} = 0{,}35\text{ mol}$.\\
Gọi $n_{\ce{Fe}} = x\text{ mol}, n_{\ce{Cu}} = y\text{ mol}$.\\
Hệ phương trình:\\
- Khối lượng: $56x + 64y = 18{,}4$ (1)\\
- Bảo toàn electron: $3x + 2y = 2 \times n_{\ce{SO2}} = 2 \times 0{,}35 = 0{,}7$ (2)\\
Giải hệ (1) và (2) ta được: $x = 0{,}1\text{ mol}; y = 0{,}2\text{ mol}$.\\
$\implies m_{\ce{Fe}} = 0{,}1 \times 56 = 5{,}6\text{ gam} \implies \%m_{\ce{Fe}} = \frac{5{,}6}{18{,}4} \times 100\% = 30{,}43\%$.\\
c) Sai vì Al và Fe bị thụ động hóa trong $\ce{HNO3}$ đặc, nguội.\\
d) Đúng vì $\ce{H2SO4}$ tan vô hạn trong nước và tỏa một lượng nhiệt rất lớn do tạo các hydrate bền.}
\end{ex}

\bigskip
\noindent
\textbf{PHẦN III (2,0 điểm): Câu trắc nghiệm yêu cầu trả lời ngắn.}\\
\textit{Thí sinh trả lời từ câu 1 đến câu 8.}

\begin{ex}
Cho các phản ứng:\\
(a) $\ce{S + O2 ->[t^\circ] SO2}$\\
(b) $\ce{S + 3F2 -> SF6}$\\
(c) $\ce{Hg + S -> HgS}$\\
(d) $\ce{H2 + S ->[t^\circ] H2S}$\\
Số phản ứng trong đó sulfur đơn chất đóng vai trò chất khử là bao nhiêu?
\shortans{2}
\loigiai{
Trong (a): $\overset{0}{\ce{S}} \rightarrow \overset{+4}{\ce{S}} + 4e \implies \ce{S}$ là chất khử.\\
Trong (b): $\overset{0}{\ce{S}} \rightarrow \overset{+6}{\ce{S}} + 6e \implies \ce{S}$ là chất khử.\\
Trong (c): $\overset{0}{\ce{S}} + 2e \rightarrow \overset{-2}{\ce{S}} \implies \ce{S}$ là chất oxi hóa.\\
Trong (d): $\overset{0}{\ce{S}} + 2e \rightarrow \overset{-2}{\ce{S}} \implies \ce{S}$ là chất oxi hóa.\\
Vậy có 2 phản ứng (a) và (b) trong đó S đóng vai trò chất khử.\\
\textbf{Đáp án: 2.}}
\end{ex}

\begin{ex}
Cho các chất: (1) $\ce{KNO3}$, (2) $\ce{HCl}$, (3) $\ce{HClO}$, (4) $\ce{H2S}$, (5) $\ce{Cu(OH)2}$, (6) $\ce{Ba(OH)2}$, (7) $\ce{NaHCO3}$. Có bao nhiêu chất điện li mạnh trong các chất trên?
\shortans{4}
\loigiai{
Chất điện li mạnh gồm muối tan, acid mạnh, base mạnh:\\
- (1) $\ce{KNO3}$ (muối tan),\\
- (2) $\ce{HCl}$ (acid mạnh),\\
- (6) $\ce{Ba(OH)2}$ (base mạnh),\\
- (7) $\ce{NaHCO3}$ (muối tan).\\
Các chất còn lại: (3) $\ce{HClO}$, (4) $\ce{H2S}$ là acid yếu; (5) $\ce{Cu(OH)2}$ là base không tan (chất điện li yếu).\\
Vậy có 4 chất điện li mạnh.\\
\textbf{Đáp án: 4.}}
\end{ex}

\begin{ex}
Cho các phản ứng sau:\\
(1) $\ce{NaOH + HCl -> NaCl + H2O}$\\
(2) $\ce{H2 + I2 <=> 2HI}$\\
(3) $\ce{CaCO3 <=> CaO + CO2}$ (trong bình kín)\\
(4) $\ce{2KClO3 ->[t^\circ] 2KCl + 3O2}$\\
Có bao nhiêu phản ứng thuận nghịch trong các phản ứng trên?
\shortans{2}
\loigiai{
Phản ứng thuận nghịch được biểu diễn bằng mũi tên hai chiều ($\rightleftharpoons$): (2) và (3).\\
Phản ứng (1) và (4) là phản ứng một chiều ($\rightarrow$).\\
Vậy có 2 phản ứng thuận nghịch.\\
\textbf{Đáp án: 2.}}
\end{ex}

\begin{ex}
Cho các nhận định sau về cấu tạo phân tử nitric acid:\\
(1) Liên kết $\ce{O - H}$ phân cực về oxygen.\\
(2) Nguyên tử N có số oxi hóa là $+5$.\\
(3) Nguyên tử N có hóa trị bằng 4.\\
(4) Phân tử $\ce{HNO3}$ chứa liên kết cho -- nhận $\ce{N -> O}$.\\
Có bao nhiêu phát biểu đúng?
\shortans{4}
\loigiai{
Cả 4 phát biểu đều đúng:\\
- (1) Đúng vì độ âm điện O ($3{,}44$) lớn hơn H ($2{,}20$) nên liên kết phân cực về O.\\
- (2) Đúng vì trong $\ce{HNO3}$, H có số oxi hóa $+1$, O có số oxi hóa $-2 \implies$ N có số oxi hóa $+5$.\\
- (3) Đúng vì N ở chu kì 2 có tối đa 4 orbital hóa trị nên hóa trị tối đa bằng 4 (IV).\\
- (4) Đúng vì trong công thức cấu tạo, liên kết giữa N và một nguyên tử O là liên kết phối trí (cho -- nhận $\ce{N -> O}$).\\
\textbf{Đáp án: 4.}}
\end{ex}

\begin{ex}
Cho phản ứng: $a\ce{Fe} + b\ce{HNO3 ->} c\ce{Fe(NO3)3} + d\ce{NO ^} + e\ce{H2O}$. Hệ số tỉ lượng $a, b, c, d, e$ là những số nguyên dương có tỉ lệ tối giản. Tổng $(a + b)$ bằng bao nhiêu?
\shortans{5}
\loigiai{
Cân bằng phản ứng oxi hóa -- khử:\\
$1\ce{Fe} + 4\ce{HNO3 ->} 1\ce{Fe(NO3)3} + 1\ce{NO ^} + 2\ce{H2O}$.\\
Các hệ số tối giản: $a = 1, b = 4, c = 1, d = 1, e = 2$.\\
Tổng $a + b = 1 + 4 = 5$.\\
\textbf{Đáp án: 5.}}
\end{ex}

\begin{ex}
Cho sơ đồ chuyển hóa nitrogen trong khí quyển thành phân đạm:\\
$\ce{N2 ->[(1)][+O2] NO ->[(2)][+O2] NO2 ->[(3)][+O2, +H2O] HNO3 ->[(4)] NO3-}$.\\
Có bao nhiêu phản ứng thuộc loại oxi hóa -- khử trong sơ đồ trên?
\shortans{3}
\loigiai{
- Phản ứng (1): $\overset{0}{\ce{N}}_2 + \overset{0}{\ce{O}}_2 \xrightarrow{3000^\circ\text{C}} 2\overset{+2}{\ce{N}}\overset{-2}{\ce{O}}$ (oxi hóa -- khử).\\
- Phản ứng (2): $2\overset{+2}{\ce{N}}\ce{O} + \overset{0}{\ce{O}}_2 \rightarrow 2\overset{+4}{\ce{N}}\ce{O}_2$ (oxi hóa -- khử).\\
- Phản ứng (3): $4\overset{+4}{\ce{N}}\ce{O}_2 + \overset{0}{\ce{O}}_2 + 2\ce{H2O} \rightarrow 4\ce{H}\overset{+5}{\ce{N}}\ce{O}_3$ (oxi hóa -- khử).\\
- Phản ứng (4): $\ce{HNO3 -> H+ + NO3-}$ (sự phân li điện li tạo ion hoặc phản ứng trung hòa tạo muối), số oxi hóa của N trước và sau đều là $+5$ (không phải phản ứng oxi hóa -- khử).\\
Vậy có 3 phản ứng oxi hóa -- khử: (1), (2), (3).\\
\textbf{Đáp án: 3.}}
\end{ex}

\begin{ex}
Cho $m$ gam Fe tác dụng hết với dung dịch $\ce{HNO3}$ loãng (dư), sau phản ứng thu được $4{,}958\text{ lít}$ (đkc) khí NO là sản phẩm khử duy nhất. Tính giá trị của $m$?
\shortans{11{,}2}
\loigiai{
Số mol khí $\ce{NO} = \frac{4{,}958}{24{,}79} = 0{,}2\text{ mol}$.\\
Phương trình hóa học: $\ce{Fe + 4HNO3 -> Fe(NO3)3 + NO ^ + 2H2O}$.\\
Theo phương trình hoặc bảo toàn electron: $n_{\ce{Fe}} = n_{\ce{NO}} = 0{,}2\text{ mol}$.\\
Khối lượng sắt: $m = 0{,}2 \times 56 = 11{,}2\text{ gam}$.\\
\textbf{Đáp án: 11,2.}}
\end{ex}

\begin{ex}
Có các thí nghiệm sau:\\
(a) Nhúng thanh sắt (iron) vào dung dịch $\ce{H2SO4}$ loãng, nguội.\\
(b) Sục khí $\ce{SO2}$ vào nước bromine.\\
(c) Cho dung dịch $\ce{BaCl2}$ tác dụng với dung dịch $\ce{H2SO4}$ loãng.\\
(d) Nhúng lá nhôm (aluminium) vào dung dịch $\ce{H2SO4}$ đặc, nguội.\\
Số thí nghiệm xảy ra phản ứng hoá học là bao nhiêu?
\shortans{3}
\loigiai{
- (a) $\ce{Fe + H2SO4 -> FeSO4 + H2 ^}$ (có phản ứng).\\
- (b) $\ce{SO2 + Br2 + 2H2O -> 2HBr + H2SO4}$ (có phản ứng làm mất màu nước bromine).\\
- (c) $\ce{BaCl2 + H2SO4 -> BaSO4 v + 2HCl}$ (có phản ứng tạo kết tủa trắng).\\
- (d) Al bị thụ động hóa trong $\ce{H2SO4}$ đặc, nguội (không xảy ra phản ứng).\\
Vậy có 3 thí nghiệm xảy ra phản ứng hóa học: (a), (b), (c).\\
\textbf{Đáp án: 3.}}
\end{ex}

\bigskip
\noindent
\textbf{\large B. PHẦN TỰ LUẬN (3,0 điểm)}

\begin{ex}
\textbf{(0,5 điểm)}
\begin{enumerate}[a)]
	\item Ở các vùng quê, người dân thường dùng phèn chua ($\ce{K2SO4.Al2(SO4)3.24H2O}$) để làm trong nước nhờ ứng dụng của phản ứng thuỷ phân ion $\ce{Al^3+}$. Giải thích?
	\item Chất hay ion nào là acid, là base trong phản ứng thuỷ phân $\ce{Al^3+}$?
\end{enumerate}
\loigiai{
a) Khi hòa tan phèn chua vào nước, ion $\ce{Al^3+}$ bị thủy phân theo cân bằng:\\
$\ce{Al^3+ + 3H2O <=> Al(OH)3 v + 3H+}$ (hoặc $[\ce{Al(H2O)6}]^{3+} + \ce{H2O <=> [Al(OH)(H2O)5]^{2+} + H3O+}$).\\
Kết tủa keo $\ce{Al(OH)3}$ dạng xốp có diện tích bề mặt lớn và khả năng hấp phụ rất cao, lơ lửng trong nước sẽ kéo dính các hạt bụi đất, cặn bẩn nhỏ li ti lắng nhanh xuống đáy làm cho nước trở nên trong sạch.\\
b) Theo thuyết Brønsted -- Lowry trong cân bằng thủy phân trên:\\
- Ion $\ce{Al^3+}$ (dạng $[\ce{Al(H2O)6}]^{3+}$) nhường proton $\ce{H+}$ nên đóng vai trò là \textbf{acid}.\\
- Phân tử $\ce{H2O}$ nhận proton $\ce{H+}$ nên đóng vai trò là \textbf{base}.}
\end{ex}

\begin{ex}
\textbf{(1,5 điểm)}
\begin{enumerate}[a)]
	\item Viết sơ đồ sản xuất nitric acid trong công nghiệp.
	\item Tính số kg dung dịch $\ce{HNO3}\ 60\%$ điều chế được từ $340\text{ kg}$ ammonia, biết rằng hiệu suất của toàn bộ quá trình là $90\%$.
\end{enumerate}
\loigiai{
a) Sơ đồ sản xuất nitric acid trong công nghiệp (phương pháp Ostwald):\\
\[ \ce{NH3 ->[+O2, \text{Pt}][850 - 900^\circ\text{C}] NO ->[+O2] NO2 ->[+O2, +H2O] HNO3} \]
Các phương trình hóa học xảy ra:\\
(1) $\ce{4NH3 + 5O2 ->[850 - 900^\circ\text{C}][\text{Pt}] 4NO + 6H2O}$\\
(2) $\ce{2NO + O2 -> 2NO2}$\\
(3) $\ce{4NO2 + O2 + 2H2O -> 4HNO3}$\\
b) Tính toán:\\
- Số mol ammonia: $n_{\ce{NH3}} = \frac{340\text{ kg}}{17\text{ kg/kmol}} = 20\text{ kmol}$.\\
- Theo sơ đồ phản ứng lí thuyết: $\ce{NH3 -> HNO3} \implies n_{\ce{HNO3 (LT)}} = n_{\ce{NH3}} = 20\text{ kmol}$.\\
- Do hiệu suất cả quá trình $H = 90\% \implies n_{\ce{HNO3 (TT)}} = 20 \times \frac{90}{100} = 18\text{ kmol}$.\\
- Khối lượng chất tan $\ce{HNO3}$ nguyên chất: $m_{\ce{HNO3}} = 18\text{ kmol} \times 63\text{ kg/kmol} = 1134\text{ kg}$.\\
- Khối lượng dung dịch $\ce{HNO3}\ 60\%$ thu được:\\
\[ m_{\text{dd }\ce{HNO3}\ 60\%} = \frac{1134 \times 100}{60} = 1890\text{ kg}. \]
\textbf{Đáp số: $1890\text{ kg}$.}}
\end{ex}

\begin{ex}
\textbf{(1,0 điểm)}
\begin{enumerate}[a)]
	\item Bình đựng dung dịch $\ce{H2SO4}$ đặc để trong không khí ẩm lâu ngày thì khối lượng bình có thay đổi không? Vì sao?
	\item Sulfur và quặng pyrite sắt là các nguyên liệu chính trong công nghiệp sản xuất sulfuric acid. Tại một nhà máy, cứ đốt cháy 1 tấn quặng pyrite sắt (chứa $84\%$ khối lượng $\ce{FeS2}$) bằng không khí, thu được tối đa $V\text{ m}^3$ khí $\ce{SO2}$ (đkc). Tính giá trị của $V$?
\end{enumerate}
\loigiai{
a) Khối lượng bình đựng dung dịch $\ce{H2SO4}$ đặc để lâu ngày trong không khí ẩm sẽ \textbf{tăng lên}.\\
\textit{Giải thích:} Do sulfuric acid đặc có tính háo nước (hút ẩm) rất mạnh, nó hút hơi nước có trong không khí ẩm xung quanh vào bình làm tăng khối lượng chất lỏng trong bình, dẫn đến khối lượng của bình tăng lên.\\
b) Tính toán:\\
- Khối lượng $\ce{FeS2}$ nguyên chất có trong 1 tấn quặng pyrite:\\
$m_{\ce{FeS2}} = 1\text{ tấn} \times 84\% = 0{,}84\text{ tấn} = 840\text{ kg} = 840\,000\text{ gam}$.\\
- Khối lượng mol của $\ce{FeS2}$: $M = 56 + 32 \times 2 = 120\text{ g/mol}$.\\
- Số mol của $\ce{FeS2}$: $n_{\ce{FeS2}} = \frac{840\,000}{120} = 7\,000\text{ mol} = 7\text{ kmol}$.\\
- Phương trình hóa học đốt cháy quặng pyrite:\\
\[ \ce{4FeS2 + 11O2 ->[t^\circ] 2Fe2O3 + 8SO2} \]
- Theo phương trình phản ứng (hoặc bảo toàn nguyên tố S):\\
$n_{\ce{SO2}} = 2 \times n_{\ce{FeS2}} = 2 \times 7\text{ kmol} = 14\text{ kmol} = 14\,000\text{ mol}$.\\
- Thể tích khí $\ce{SO2}$ thu được tối đa ở điều kiện chuẩn (đkc, $25^\circ\text{C}, 1\text{ bar}$, $24{,}79\text{ L/mol}$):\\
\[ V = 14\,000 \times 0{,}02479\text{ m}^3 = 347{,}06\text{ m}^3. \]
\textit{(Ghi chú: Nếu tính theo đktc cũ $22{,}4\text{ L/mol}$ thì $V = 14\,000 \times 0{,}0224 = 313{,}6\text{ m}^3$).}\\
\textbf{Đáp số: $V = 347{,}06\text{ m}^3$ (đkc).}}
\end{ex}

\bigskip
\noindent
{\color{hfRed}\textbf{\large BẢNG TỔNG HỢP ĐÁP ÁN ĐỀ SỐ 102}}

\smallskip
\noindent
\textbf{1. PHẦN I (Trắc nghiệm nhiều phương án):}
\begin{center}
\begin{tabular}{|c|c|c|c|c|c|c|c|c|c|c|c|}
\hline
\textbf{1. A} & \textbf{2. A} & \textbf{3. D} & \textbf{4. B} & \textbf{5. D} & \textbf{6. B} & \textbf{7. B} & \textbf{8. D} & \textbf{9. A} & \textbf{10. A} & \textbf{11. D} & \textbf{12. B} \\ \hline
\end{tabular}
\end{center}

\noindent
\textbf{2. PHẦN II (Trắc nghiệm đúng sai):}
\begin{itemize}
	\item \textbf{Câu 1:} a) Đúng \quad|\quad b) Sai \quad|\quad c) Sai \quad|\quad d) Sai
	\item \textbf{Câu 2:} a) Sai \quad|\quad b) Đúng \quad|\quad c) Sai \quad|\quad d) Đúng
\end{itemize}

\noindent
\textbf{3. PHẦN III (Trả lời ngắn):}
\begin{center}
\begin{tabular}{|c|c|c|c|c|c|c|c|}
\hline
\textbf{Câu 1} & \textbf{Câu 2} & \textbf{Câu 3} & \textbf{Câu 4} & \textbf{Câu 5} & \textbf{Câu 6} & \textbf{Câu 7} & \textbf{Câu 8} \\ \hline
2 & 4 & 2 & 4 & 5 & 3 & 11,2 & 3 \\ \hline
\end{tabular}
\end{center}

\begin{center}
\textbf{--------- HẾT ---------}\\
\textit{(Học sinh không được sử dụng tài liệu. Giám thị không giải thích gì thêm)}
\end{center}

\end{document}
